% 编译方式：使用支持中文的 XeLaTeX 连续编译两次。
% 排版依据：旧版论文草稿的文档类、页面设置、标题层级、编号和三线表结构。
% 当前为阶段稿：摘要仅覆盖问题一；问题二至四已完成基准求解与结果分析。
% 问题二至四的方法章节仍需按实际代码统一，敏感性结果不得超出口径说明。
% 后续正文在参考文献之前续写，保留现有格式；全题完成后统一修订题目和摘要。
\documentclass[UTF8,a4paper,12pt,fontset=fandol]{ctexart}
\usepackage[left=2.5cm,right=2.5cm,top=2.4cm,bottom=2.4cm]{geometry}
\usepackage{amsmath,amssymb,bm}
\usepackage{booktabs,longtable,tabularx,array,multirow}
\usepackage{graphicx,float,subcaption}
\usepackage{xcolor}
\usepackage{enumitem}
\usepackage{setspace}
\usepackage{fancyhdr}
\usepackage{hyperref}
\usepackage{listings}
\usepackage{caption}
\usepackage{chngcntr}

\hypersetup{colorlinks=true,linkcolor=black,citecolor=black,urlcolor=blue}
\graphicspath{{figures/}{../figures/}}
\captionsetup{font=small,labelsep=quad}
\captionsetup[table]{position=top}
\captionsetup[figure]{position=bottom}
\setstretch{1.22}
\setlength{\parindent}{2em}
\setlength{\parskip}{0pt}
\setlist{nosep,leftmargin=2em}
\pagestyle{fancy}
\fancyhf{}
\cfoot{\thepage}
\renewcommand{\headrulewidth}{0pt}
\renewcommand{\arraystretch}{1.18}
\newcolumntype{Y}{>{\centering\arraybackslash}X}
\newcolumntype{L}{>{\raggedright\arraybackslash}X}
\newcommand{\dif}{\mathop{}\!\mathrm{d}}
\newcommand{\VaR}{\operatorname{VaR}}
\newcommand{\CVaR}{\operatorname{CVaR}}
\numberwithin{equation}{section}
\counterwithin{figure}{section}
\counterwithin{table}{section}
\renewcommand{\thefigure}{\thesection-\arabic{figure}}
\renewcommand{\thetable}{\thesection-\arabic{table}}

\lstset{
  language=Python,
  basicstyle=\ttfamily\scriptsize,
  keywordstyle=\color{blue!70!black},
  commentstyle=\color{green!45!black},
  stringstyle=\color{orange!70!black},
  numbers=left,
  numberstyle=\tiny\color{gray},
  frame=single,
  breaklines=true,
  showstringspaces=false,
  tabsize=4,
  columns=fullflexible
}

\title{\heiti\bfseries 基于两层词典序混合整数线性规划模型的微网与外部电网单日电力调控问题研究}
\author{}
\date{}

\begin{document}
\maketitle
\vspace{-2.2em}

% 摘要采用局部紧凑排版；其余正文保留旧稿的默认字号和行距。
% 摘要的一页要求需在完整编译后的页面中最终核验。
\begingroup
\small
\setstretch{1.08}
\begin{abstract}
光伏出力与居民用电负荷在时间上的不匹配，以及分时电价的差异，使微网面临购电成本控制与储能安全运行的双重要求。针对给定负荷、光伏预测及电价条件下的单日电力调控问题，本文联合优化外部电网购电量与储能充放电策略，在逐时段满足负荷需求的前提下，降低全天购电费用，并通过约束回代、基准对照和效率敏感性分析检验方案的可行性与经济性。

首先，将全天划分为144个十分钟时段，统一功率与电量的换算口径，以购电量、充电量、放电量、弃光量及储能状态为变量，建立\textbf{混合整数线性规划模型}。模型引入二元变量约束充放电互斥，综合考虑供需平衡、储能状态递推、容量边界、充放电功率上限及日初日末储能相等条件。在允许弃光、不向外部电网售电的假设下，将充电和放电单程效率均设为90\%，储能运行区间设为1200--10800千瓦时，最大充放电功率均为5000千瓦，确保调度方案满足设备运行条件。

其次，采用\textbf{两层词典序优化方法}：第一层最小化全天购电费用；第二层在购电费用不超过第一层最优值加给定容差的条件下，最小化储能充放电总吞吐量，避免人为设定折旧权重改变首要目标。模型通过低价时段购电充能、光伏富余时段吸纳余电及高价时段放电，实现跨时段能量配置，并生成逐时段购电方案与分时段充放电汇总。

求解结果表明，优化后全天购电费用为\textbf{35126.95元}，较无储能基准48052.05元降低12925.10元，\textbf{降幅为26.90\%}；全天购电量为59482.70千瓦时，充电量与放电量分别为20740.67千瓦时和16799.94千瓦时。在给定光伏预测情景下，全天弃光量为零，日末储能恢复至6000千瓦时，全部时段均满足负荷需求，且不存在同时充放电现象。

最后，对输出方案进行数值检验，供需平衡与储能递推的最大残差分别约为$2.27\times10^{-13}$和$8.31\times10^{-13}$千瓦时，储能状态全程满足上下界约束。效率敏感性分析显示，单程效率由80\%提高至95\%时，最优购电费用由38223.65元降至33767.03元；若将题设90\%解释为往返效率，费用为33801.50元，表明效率口径会影响经济性评估。研究表明，该模型能够在满足储能物理约束的条件下有效降低购电成本、促进光伏消纳，具有结构清晰、结果可解释及便于复现的优势，可为后续考虑预测误差与波动电价的调度研究提供基础。

\par\medskip
\noindent\textbf{关键词：}混合整数线性规划；两层词典序优化；微网电力调控；储能调度；光伏消纳
\end{abstract}
\endgroup
\clearpage

\section{问题重述}
\label{sec:restatement}

\subsection{问题背景}

含光伏与储能的并网型微网能够在本地协调发电、用电及能量存储，但其经济运行受到供需错配与预测不确定性的共同制约。一方面，光伏出力具有明显的日内变化，发电高峰未必与居民用电高峰重合，局部时段可能出现电能富余，而其他时段仍需依赖外部电网。另一方面，储能虽然能够转移电能的使用时刻，却受到容量、充放电功率和能量转换损耗的限制。因此，仅依据“低价充电、高价放电”的经验规则，难以协调光伏消纳、购电支出和后续时段的供电需求。

当调度由确定性单日计划扩展至连续运行时，问题进一步表现为预测与决策之间的相互影响。计划购电偏少会增加高价紧急购电的需求；计划偏多则可能形成无法充分利用、但仍需支付费用的购电量。即使日内能够更新光伏预报，调整计划也伴随着额外交易费用，因此，预测更新能否带来实际收益，需要通过完整的调度成本衡量。

近年来的相关研究为上述问题提供了方法依据。2022年发表于《可再生能源》的研究将实际负荷、可再生能源数据及预测误差纳入经济模型预测控制框架，强调利用滚动决策适应实际运行状态\cite{economic2022}。2023年的风险感知微网能源管理研究指出，应系统考虑发电与用电预测的不确定性，同时控制随机优化的计算负担\cite{riskaware2023}。据此，本题需要围绕“物理约束下的经济调度---预测误差下的补救决策---信息更新下的动态调整”建立统一研究框架。

\subsection{研究对象与约束条件}

研究对象为由小区负荷、光伏发电、储能设备及外部电网构成的并网型微网。调度以十分钟为基本时间间隔，每日包含144个时段，需要联合确定各时段的计划购电量和储能充放电量，并在需要时安排计划调整与紧急购电。

储能设备额定容量为12000千瓦时，允许运行电量为1200--10800千瓦时，最大充电功率和最大放电功率均为5000千瓦。2025年1月1日零时的初始储电量为6000千瓦时，题设充放电效率为90\%。其中，问题一还要求日初与日末储电量相同。

数据包括典型日的负荷、光伏预测与电价，以及2025年全年的实际负荷、实际光伏出力、多时点光伏预报和波动电价。全年策略需按要求输出2月1日至12月31日的调度记录，并展示3月20日、6月21日、9月23日和12月21日的代表性结果。预测数据与实际数据应按发布时间及目标时段分别管理，避免把事后观测值作为事前已知条件。

\subsection{核心研究任务}

\textbf{研究任务一：确定性条件下的购电与储能协同优化，对应原题问题一。}

在当日负荷、电价及光伏预测已知的条件下，建立满足供需平衡与储能运行限制的单日调度模型。重点研究储能应在何时吸收富余光伏、何时利用低价电补能，以及何时释放电能替代高价购电，使日初日末储电量相同条件下的全天购电费用最小，并形成可执行的逐时段调度方案。

\textbf{研究任务二：预测不确定性下的日前计划与日内调整，对应原题问题二、三。}

首先，在每日电价曲线相同、负荷与光伏实际出力变化的条件下，利用历史信息制定零时购电计划，并以交易时刻电价五倍的紧急购电弥补实际供电缺口。其次，引入零时、六时、十二时和十八时发布的光伏预报，在已执行决策不可更改的前提下更新剩余时段的购电与储能安排。研究重点是平衡提前采购、日内调整和紧急补购之间的成本关系，并判断增加预报时点的边际经济价值。

\textbf{研究任务三：波动电价下的策略扩展与适应性评价，对应原题问题四。}

将固定日内电价曲线替换为实际波动电价，分别研究仅进行日前计划和允许日内调整两种机制。重点分析电价不确定性与供需预测误差对储能使用时机的共同影响，在一致的数据范围、运行边界和结算口径下比较两种策略，为不同交易条件下的购电决策提供依据。

% 问题分析采用独立页面与局部紧凑排版；不混入尚未完成的求解结果。
\clearpage
\section{问题分析}
\label{sec:analysis}
\begingroup
\small
\setstretch{1.08}

\textbf{总体思路。}本题整体采用“数据处理与因果预测---约束优化---滚动执行---对照检验”的思路。以供需平衡和储能状态递推构成统一物理模型，逐步增加紧急购电、计划调整和电价不确定性。购电、调整与补救费用共同构成总成本，属于具有多个成本分项的单目标优化；问题一再以两层词典序方式处理等价方案选择。

\textbf{问题1分析。}关键矛盾是降低当前购电费用与保留后续可用储能之间的跨时段权衡。首先统一功率、电量与时间索引，识别净负荷及光伏富余区间；其次建立购电、充放电和储能状态变量，以二元变量保证充放电互斥，并施加容量、功率及首末储能约束；最后先最小化购电费，再在规定成本容差内最小化吞吐量，通过约束回代和线性松弛下界检验方案。

\textbf{问题2分析。}关键矛盾是计划不足引发高价补购，而过度采购仍需承担计划费用。先按时间顺序利用历史数据预测负荷和光伏，再根据预测残差构造保留时间相关性的情景；随后建立日前计划与实时补救相衔接的两阶段模型。执行时仅依据当时可用信息调整储能和紧急购电，保持跨日状态连续，避免利用未来实际值获得不可实现的策略优势。

\textbf{问题3分析。}关键矛盾是新预报可能降低供需偏差，但修改计划需要支付费用。先将整点预报统一至十分钟网格，再于各发布时间更新剩余时域的预测，固定已执行决策，以当前储能和有效购电计划为起点进行滚动优化。结算时区分原始计划、增减调整量和紧急购电。新增预报时点通过成对回测比较净成本变化，其收益须覆盖信息获取及调整成本。

\textbf{问题4分析。}关键矛盾是储能同时承担电价套利与供电缓冲功能，单纯追逐价差可能削弱应对缺口的能力。先明确未来电价的可获得范围，对未知电价建立预测及联合误差情景；再分别嵌入问题二、三的交易结构，保持一致的评价条件。最后结合总成本、紧急购电、尾部成本及参数敏感性，分析策略对价格波动的适应能力。

% 以下流程框架默认不排印，避免与分析正文重复并占用单页篇幅。
% 如后续需要流程图，可据此绘制并通过图环境插入。
% 问题一：数据与单位核验→确定性混合整数优化→词典序择优→约束及下界检验。
% 问题二：历史信息→预测与误差情景→日前计划→实时补救→跨日状态更新。
% 问题三：新预报到达→固定已执行决策→剩余时域优化→费用结算→新增预报价值评价。
% 问题四：电价信息边界判断→联合预测与情景生成→分别运行两种交易机制→统一对照评价。
\par
\endgroup
\clearpage

\section{模型假设}
\label{sec:assumptions}

\begin{enumerate}
  \item \textbf{各时段数据可用其平均水平表示，且调度不得使用未来信息。}将连续运行过程划分为十分钟时段，认为时段内负荷、光伏功率和电价保持不变；问题二至四仅使用决策时点前已观测或已发布的数据，整点预报按统一规则映射到十分钟网格，据此建立可实际执行的预测与滚动优化模型。
  \item \textbf{储能效率和可用容量在研究期内保持不变。}将题目给出的90\%解释为充电、放电各自的单程效率，忽略自放电以及温度、功率和循环老化对效率与容量的影响，据此采用线性的储能状态递推关系。
  \item \textbf{微网可等效为单节点系统，富余电量允许舍弃但不能出售。}忽略内部线路损耗、电压和无功约束，认为外部电网在题设情景下能够正常供电；光伏富余量可以弃置，已购但未被负荷或储能利用的电量不退款、不产生售电收益，据此闭合各时段的能量平衡。
  \item \textbf{计划购电与调整费用按不重复计价原则结算。}保留零时原计划购电费，上调电量按交易时刻电价的1.5倍加收，下调电量按0.5倍计收违约费用，紧急购电按5倍电价另行结算；同一时段经过多次滚动修正时，只按最终生效计划相对原计划的增减量计费。
  \item \textbf{相邻日期的储能状态连续，基准方案要求每日末储电量回到当日初值。}问题一依题意满足日初、日末储电量相等；问题二至四以前一日实际末状态作为次日初状态，并采用每日首末相等的统一边界，以避免通过期末放空储能虚降成本，保证不同策略具有相同的比较基础。
\end{enumerate}

\section{符号说明}
\label{sec:notation}

表\ref{tab:notation}给出全题共用及后续模型拟采用的主要符号。功率与电量、交流侧电量与电池内部电量以及预测值与实际执行值分别表示，避免不同建模阶段出现口径混用。

\begingroup
\small
\setlength{\tabcolsep}{4pt}
\begin{longtable}{>{\raggedright\arraybackslash}p{0.25\textwidth}>{\raggedright\arraybackslash}p{0.53\textwidth}>{\raggedright\arraybackslash}p{0.15\textwidth}}
\caption{主要符号及其含义}\label{tab:notation}\\
\toprule
符号 & 含义 & 单位\\
\midrule
\endfirsthead
\multicolumn{3}{c}{表\thetable\quad 主要符号及其含义（续）}\\
\toprule
符号 & 含义 & 单位\\
\midrule
\endhead
\midrule
\multicolumn{3}{r}{续下页}\\
\endfoot
\bottomrule
\endlastfoot
$t,\mathcal T$ & 十分钟时段索引及所研究时段集合 & 无量纲\\
$k,\tau_k$ & 滚动决策序号及对应的决策时点 & 无量纲；时刻\\
$\omega,\Omega$ & 不确定性情景索引及情景集合 & 无量纲\\
$\pi_\omega$ & 情景$\omega$的发生概率 & 无量纲\\
$\Delta t$ & 调度时间间隔，取$1/6$小时 & 小时\\
$\mathcal I_k$ & 决策时点$\tau_k$已获得的观测与预报信息集合 & ---\\
$P_t^L,P_t^{PV}$ & 时段平均负荷功率、可用光伏功率 & 千瓦\\
$L_t,R_t$ & 时段负荷电量、可用光伏电量 & 千瓦时\\
$\widehat L_{t\mid k},\widehat R_{t\mid k}$ & 在决策时点$\tau_k$对时段$t$的负荷、光伏电量预测 & 千瓦时\\
$N_t$ & 净负荷电量，定义为$L_t-R_t$ & 千瓦时\\
$p_t,\widehat p_{t\mid k}$ & 时段结算电价及在决策时点形成的对应电价预测 & 元／千瓦时\\
$g_t^0$ & 每日零时确定的时段原始计划购电量 & 千瓦时\\
$g_{t\mid k}$ & 第$k$次决策形成的时段非紧急购电计划量 & 千瓦时\\
$g_t$ & 时段最终生效的非紧急购电量 & 千瓦时\\
$\delta_t^+,\delta_t^-$ & 最终生效购电量相对零时原计划的上调量、下调量 & 千瓦时\\
$e_t$ & 时段紧急购电量 & 千瓦时\\
$c_t,d_t$ & 时段交流侧充电量、交流侧放电量 & 千瓦时\\
$w_t,s_t$ & 时段弃光电量、已购买但未利用的电量 & 千瓦时\\
$E_t$ & 时段结束时的电池内部储电量 & 千瓦时\\
$E^{\mathrm{ini}},E^{\mathrm{ter}}$ & 所研究区间的初始、目标终端储电量 & 千瓦时\\
$E^{\min},E^{\max}$ & 允许储电量下限、上限，分别为1200、10800 & 千瓦时\\
$P^{ch,\max},P^{dis,\max}$ & 最大充电功率、最大放电功率，均为5000 & 千瓦\\
$\eta_{ch},\eta_{dis}$ & 充电效率、放电效率 & 无量纲\\
$u_t^{ch},u_t^{dis}$ & 时段充电、放电状态二元变量 & 无量纲\\
\shortstack[l]{$C_{\mathrm{plan}},C_{\mathrm{adj}}$\\$C_{\mathrm{emg}}$} & 评价区间内计划购电费、调整费、紧急购电费 & 元\\
$C_{\mathrm{tot}}$ & 评价区间内实际总购电费用 & 元\\
$C^*,\varepsilon_C$ & 问题一第一层最优购电费、第二层允许的费用容差 & 元\\
$Q$ & 评价区间内储能交流侧充放电总吞吐量 & 千瓦时\\
$V_\tau$ & 增加预报时点$\tau$所带来的预期净成本节约 & 元\\
\end{longtable}
\endgroup

为保持后续公式一致，统一以下口径：
\begin{enumerate}
\item \textbf{功率与电量严格区分。}$L_t=P_t^L\Delta t$、$R_t=P_t^{PV}\Delta t$。购电及充放电变量均表示时段电量，计算费用时不再重复乘以时间间隔。
\item \textbf{储能侧与交流侧严格区分。}$c_t,d_t$为交流侧电量，$E_t$为电池内部电量，效率仅在储能状态递推中计入。
\item \textbf{计划与执行严格区分。}未调整时$g_t=g_t^0$；允许调整时$g_t=g_t^0+\delta_t^+-\delta_t^-$。两类调整量分别表示同一差额的正、负部分，不应同时为正。
\item \textbf{预测与情景保留索引。}帽号表示预测，竖线后的$k$表示预测形成时点；情景模型中的变量按需增加$\omega$索引，实际执行结果不得混用情景预测值。
\end{enumerate}

\section{问题二至问题四结果分析}
\label{sec:q234-results}

\subsection{统计口径与总体结果}

问题二至问题四均以2025年1月作为预测历史热启动期，正式评价区间为2025年2月1日至12月31日，共334天、每种策略48096个十分钟时段。该区间实际负荷总量为3700.8079万千瓦时，实际光伏发电量为1906.8890万千瓦时，实际净负荷为1793.9189万千瓦时。四种策略的核心结果见表\ref{tab:q234-total}。

\begin{table}[H]
\centering
\caption{问题二至问题四正式评价期总体结果}
\label{tab:q234-total}
\resizebox{\textwidth}{!}{%
\begin{tabular}{lrrrrr}
\toprule
策略 & 总成本/万元 & 计划购电/万千瓦时 & 调整后购电/万千瓦时 & 紧急购电/万千瓦时 & 剩余电量/万千瓦时\\
\midrule
问题二：固定价、日前计划 & 1720.9391 & 2277.7914 & 2277.7914 & 83.0700 & 432.6786\\
问题三：固定价、四时点滚动 & 1611.7656 & 2040.4779 & 2027.7637 & 87.9152 & 193.9985\\
问题四之日前策略：波动价 & 1805.3451 & 2280.9361 & 2280.9361 & 83.0700 & 432.6786\\
问题四之滚动策略：波动价 & 1685.0318 & 2043.6907 & 2030.2772 & 87.9308 & 193.4288\\
\bottomrule
\end{tabular}}
\end{table}

四个方案的最大供需平衡残差均为$2.274\times10^{-13}$千瓦时，最大SOC递推残差均为$9.095\times10^{-13}$千瓦时，成本对账误差为0，同时充放电时段数为0。所有时段SOC均处于1200---10800千瓦时，单时段充、放电量均未超过833.3333千瓦时，对应功率上限5000千瓦。因此，以下经济分析建立在物理可行且数值回代通过的调度结果上。

\subsection{问题二：固定电价下的日前计划}

\subsubsection{基础分析}

\textbf{数值解读。}正式评价期内，模型计划购电2277.7914万千瓦时、紧急购电83.0700万千瓦时，总购电费用为1720.9391万元。其中计划购电费1396.2996万元、紧急购电费324.6395万元。紧急电量仅占实际外购电量的3.52\%，却贡献18.86\%的总成本，说明五倍应急电价使少量预测缺口成为费用的重要来源。

模型产生432.6786万千瓦时剩余电量，占计划购电量约19.00\%。全期充电706.6517万千瓦时、放电572.3879万千瓦时，两者相差134.2638万千瓦时。计划购电与紧急购电合计超过实际净负荷的部分，可由剩余电量与储能效率损失共同解释，因此不是能量平衡错误，而是保守预测和储能损耗付出的代价。

\textbf{统计描述。}日总成本均值为51525.12元，样本方差为$3.3828\times10^8$元$^2$，标准差为18392.52元，变异系数为0.357；中位数为53862.89元，取值范围为17381.32---105636.95元，95\%分位数为83212.88元。最高成本出现在6月11日，最低成本出现在4月26日，说明季节性净负荷和预测误差使日成本差异较大。

日紧急购电量均值为2487.12千瓦时、标准差为3408.66千瓦时、变异系数为1.371，中位数仅799.04千瓦时，而最大值为15976.49千瓦时，呈明显右偏分布。334天中有331天发生紧急购电，共涉及13720个时段，占28.53\%；最大单时段紧急购电量为435.10千瓦时，发生于6月11日14:10。6月和12月紧急购电合计42.8887万千瓦时，占全期紧急电量51.63\%，风险具有季节集中性。

储能在11970个时段充电、12402个时段放电，SOC均值为6941.13千瓦时、标准差为3502.11千瓦时，共有4367个时段触及下界、7646个时段触及上界。按额定容量估算的等效放电循环约477次，若计入电池衰减，当前费用结果可能偏乐观。

\subsubsection{深层分析}

\textbf{关联性与约束匹配。}日紧急购电量与日总成本的相关系数为0.636，证实高额应急电价是成本波动的重要驱动因素。时段购电量与电价的相关系数为$-0.522$，说明储能能够把普通购电向低价时段转移。日剩余电量与日成本的相关系数为$-0.619$，这是由于光伏充足的低成本日期往往同时出现较多富余，不能把该相关关系解释为增加剩余电量会降低成本。

方案满足五倍紧急购电、SOC安全区间、功率上限、能量平衡和充放电互斥等要求。当前程序还要求每日结束SOC回到当日开始值，这比单纯保持跨日连续更保守，能够避免期末放空电池，但会限制跨日套利。

\textbf{敏感性。}将负荷分位数由0.80降至0.72、光伏分位数由0.20升至0.22后，策略更激进，紧急购电量增至99.3931万千瓦时，总成本升至1752.6785万元，较基准增加1.84\%。反之，将负荷分位数升至0.88、光伏分位数降至0.18后，计划购电量增加2.21\%，紧急购电量下降17.99\%，总成本降至1698.0183万元，较基准降低1.33\%。这说明在五倍应急电价下，当前0.80/0.20并非所测试网格中的最优参数，正式定稿前应扩大参数网格并使用独立验证集选参。

单程效率由0.90降低10\%至0.81时，总成本增加5.57\%；提高10\%至0.99时，总成本下降5.34\%。在保持既有调度不变的条件下，应急电价由5倍变为4.5倍或5.5倍，总成本分别变化$-1.89\%$和$+1.89\%$。后一结果只反映结算敏感性，若倍率事前改变，还应重新选择预测分位数和购电计划。

\textbf{实际意义。}问题二的策略本质是以适度增加普通购电和剩余电量为代价，降低高价缺口风险。实际部署应重点改善6月、12月的预测和备用配置，并将预测分位数与应急价格、剩余电量处置成本和电池衰减成本联合校准，而不能只追求预测误差最小。

\subsection{问题三：固定电价下的四时点滚动调整}

\subsubsection{基础分析}

\textbf{数值解读。}问题三在0:00形成原始计划，并于6:00、12:00、18:00依据新预报从对应整点起始区间重算未执行时段。原始计划购电2040.4779万千瓦时，最终生效购电2027.7637万千瓦时；其中上调65.9703万千瓦时、下调78.6845万千瓦时，净下调12.7142万千瓦时。由于上调按1.5倍、下调只按0.5倍冲减，虽然净调整方向向下，调整成本仍为43.4055万元。

总成本为1611.7656万元，较问题二减少109.1734万元，降幅6.34\%。剩余电量由432.6786万千瓦时降至193.9985万千瓦时，下降55.16\%；紧急购电由83.0700万千瓦时增加至87.9152万千瓦时，增加5.83\%。因此，滚动策略的主要收益不是使紧急电量最小，而是利用新信息减少过量计划和无效剩余。

\textbf{统计描述。}日总成本均值为48256.46元，样本方差为$2.5274\times10^8$元$^2$，标准差为15897.76元，中位数为50448.07元，范围为16898.58---85115.41元，95\%分位数为72798.93元。与问题二相比，其均值下降6.34\%，标准差下降13.56\%，95\%分位数下降12.51\%，最大日成本下降19.43\%，说明滚动更新同时改善平均成本和高成本尾部。

问题三每天都出现少量紧急购电，共22728个时段，占47.26\%。但是日紧急量标准差由3408.66千瓦时降至2187.65千瓦时，最大日紧急量由15976.49千瓦时降至10904.09千瓦时，最大单时段量由435.10千瓦时降至252.75千瓦时。紧急购电由较少发生但缺口较大，转为更频繁但单次较小，极端缺口受到抑制。

334天中，滚动策略有231天比日前策略成本低，占69.16\%，另有103天成本更高。逐日节省均值为3534.86元、中位数为2405.75元、标准差为7743.30元；6月11日最多节省35702.34元，7月18日反而多支出19017.93元。因此，滚动调整在全年总体有效，但并非逐日严格占优。

\subsubsection{深层分析}

\textbf{关联性与机制。}日紧急购电与日总成本的相关系数由0.636降至0.545，日成本对单一缺口的依赖有所减弱。滚动方案相对问题二减少计划购电248.9254万千瓦时、减少剩余电量239.0143万千瓦时，而紧急电量只增加3.1251万千瓦时，说明新预报的主要价值在于识别并取消原计划中的冗余购电。由于调整已计入非对称惩罚，节省额不是忽略交易成本得到的理想值。

四个代表日反映收益的非均匀性：3月20日和9月23日总成本分别由56264.07元降至52588.89元、由54167.94元降至51464.18元；但6月21日和12月21日分别由23997.63元升至29839.11元、由73568.03元升至75208.68元。因此，运营中宜设置调整后的预期节省阈值，而不应在每个发布时点机械地大幅修改计划。

\textbf{新增预报时点敏感性。}在3:00、9:00、15:00、21:00增加基于最新实测偏差的代理更新后，八时点方案总成本为1618.6479万元，比四时点方案增加6.8823万元，毛变化为$-0.43\%$，平均每天增加206.06元。逐日差额标准差为2359.83元；125天节省、209天增支。即使暂不计新增系统成本，当前代理方案也没有显示正的经济价值。

该八时点结果来自对既有预报进行实测偏差修正的代理，而非题目直接提供的新预报产品，只能作为毛信息价值估计或上界，不能直接断言必须增加四个时点。

\textbf{其他敏感性与实际意义。}单程效率降至0.81时，问题三总成本增加5.62\%；提高至0.99时下降5.30\%。应急电价倍率上下浮动10\%时，在固定调度下总成本约变化$\pm2.06\%$。滚动更新不能替代高效率储能，但能够通过减少计划冗余降低费用和高成本尾部。新增信息系统应优先改善既有整点预报质量，而不是机械增加预报频次。

\subsection{问题四：波动电价下的日前与滚动策略}

\subsubsection{基础分析}

\textbf{数值解读。}正式评价期波动电价均值为0.7575元/千瓦时，低于固定价均值0.7662元/千瓦时；但其标准差为0.3361元/千瓦时、变异系数为0.444，分别高于固定价的0.3133元/千瓦时和0.409，价格范围也由0.3713---1.3952元/千瓦时扩大至0.0076---1.7936元/千瓦时。均价略低而总成本上升，说明费用还取决于高价时段与高净负荷、预测缺口的时间重合。

波动价下，日前策略总成本为1805.3451万元，比固定价日前策略增加84.4060万元，增幅4.90\%；滚动策略总成本为1685.0318万元，比固定价滚动策略增加73.2662万元，增幅4.55\%。两种波动价策略的紧急电量分别为83.0700万千瓦时和87.9308万千瓦时，与对应固定价策略基本相同，成本增加主要来自相近购电需求在不同价格路径下的结算，而不是缺口电量明显扩大。

\textbf{统计描述。}问题四之日前策略的日成本均值为54052.25元，样本方差为$5.7844\times10^8$元$^2$，标准差为24050.74元，中位数为55402.68元，范围为12928.98---129418.96元，95\%分位数为96346.16元。滚动策略日成本均值为50450.05元，样本方差为$4.5006\times10^8$元$^2$，标准差为21214.66元，中位数为51486.76元，范围为12577.85---101744.86元，95\%分位数为85368.19元。相比固定价，对应日成本标准差分别增加30.76\%和33.44\%，价格波动明显放大经营费用的不确定性。

允许滚动调整后，波动价总成本降低120.3133万元，降幅6.66\%；334天中有214天节省、120天增支。逐日节省均值为3602.19元、中位数为2073.48元、标准差为8897.21元；12月2日最多节省44167.01元，7月18日最多增支17325.19元。滚动更新在波动价下的绝对节省比固定价多11.1399万元，说明价格波动提高了及时修正购电时点的经济价值。

\subsubsection{深层分析}

\textbf{关联性与约束匹配。}问题四之日前和滚动策略的时段购电量与电价相关系数分别为$-0.389$和$-0.364$，均保持负相关，说明储能仍遵循低价充电、高价放电的基本规律；但其绝对值小于固定价下的$-0.522$和$-0.503$，表明实时供需缺口、预测更新和终端SOC约束削弱了纯套利自由度。四个方案均频繁触及SOC上下界，储能同时承担价格套利和预测误差缓冲功能。

问题四之日前策略与问题二的紧急购电量完全一致，滚动策略只比问题三多156.69千瓦时。这说明当前实现中价格路径主要改变计划购电和储能的时点分布，能量缺口基本由负荷、光伏预测误差决定，不能据此声称波动电价提高或降低了供电可靠性。

\textbf{敏感性。}单程效率由0.90降至0.81时，波动价日前和滚动策略总成本分别增加5.31\%和5.44\%；提高至0.99时分别下降5.92\%和6.00\%，波动价环境下高效率储能的价值略高。日前策略采用较激进的0.72/0.22分位数时，总成本上升2.01\%；采用较保守的0.88/0.18分位数时，总成本下降1.51\%。应急电价倍率上下浮动10\%时，两种策略成本分别约变化$\pm1.99\%$和$\pm2.01\%$。

\textbf{实际意义与适用边界。}波动价扩大峰谷套利机会，也放大高价缺电风险。单看平均电价会低估实际费用，园区应联合考虑净负荷预测、价格路径和储能SOC，而不能把最低价充电、最高价放电作为唯一规则。四时点滚动策略在两种价格机制下均降低约7\%的总成本，方向性价值较稳定。

当前问题四代码假定当天144点电价在0:00全部可知。因此这些结果适用于日前公布完整价格路径的交易机制，也可视为该信息条件下的基准。若市场只逐时发布电价，现有结果包含价格前视信息，必须改用因果价格预测重新回测，不能把表\ref{tab:q234-total}中的金额直接视为可实现策略的确定费用。

\subsection{综合结论与结果边界}

问题二表明，高额应急电价使保守计划具有合理性，但过度计划会产生较多剩余电量；问题三表明，日内滚动更新可通过削减计划冗余使固定价总成本下降6.34\%，并降低高成本尾部；问题四表明，波动电价使总成本提高约4.5\%---4.9\%、日成本标准差提高约31\%---33\%，而滚动策略仍可节省6.66\%。

现阶段数字可作为基准方案结果，但定稿前还需处理三项边界。第一，0.88/0.18分位数组合优于当前0.80/0.20，需通过滚动训练集和验证集系统选参后决定是否重跑正式结果。第二，程序采用每日末SOC回到日初SOC的保守条件，模型章节必须保持同一口径。第三，当前未计电池衰减，而各方案按额定容量估算约453---488次等效放电循环，若循环成本较高，经济性结论需要重新评估。这些限制不改变滚动信息总体有价值、价格波动放大费用风险的方向性结论，但会影响精确金额和最优参数。

% ==================== 后续正文接续位置 ====================
% 用户要求后续步骤继续以本文件格式输出并更新本稿。
% 下一阶段新增章节放在此处，参考文献始终置于完整正文之后。
% 正文一级标题使用 section，二级标题使用 subsection，三级标题使用 subsubsection。
% 公式使用 equation、align 等环境；表格使用三线表；全局符号与上表同步。
% 未完成求解的内容不得写成已完成结果；新增结论必须能够追溯到结果文件。

\clearpage
\begin{thebibliography}{99}
% 以下保留前述已核验的题名中文译名、年份、卷页与原文链接。
% 正式定稿前补齐作者和原文题名，并按最终要求统一参考文献著录格式。
\bibitem{economic2022}
基于实际需求、可再生能源数据及预测误差的孤岛微网经济预测控制（题名中文译名）.
可再生能源, 2022, 194: 647--658.
\href{https://doi.org/10.1016/j.renene.2022.05.103}{原文链接}.

\bibitem{riskaware2023}
含电池储能与可再生能源的微网风险感知随机能源管理（题名中文译名）.
国际自动控制联合会论文集, 2023, 56(2): 8445--8450.
\href{https://doi.org/10.1016/j.ifacol.2023.10.1042}{原文链接}.
\end{thebibliography}

\end{document}
